% GENERATED FILE. Edit canonical lessons, then run npm run generate:books.
% canonical-source-hash: fnv1a64:75e543f1447c3fd4
% canonical-lessons: ZH-C449-yeshi, ZH-C449-kouwei, ZH-W449-wei-taste, ZH-C449-fengwei, ZH-C449-qingdan, ZH-C449-youni

\chapter{Tastes}
\label{ch:zh-tastes}
\hlchaptermodality{449}

\begin{chapteropening}
\textbf{By the end of this chapter:} \emph{I can talk about a night market, tastes in food, local flavours, light food and greasy food.}

Complete the last lesson of chapter 449: I can talk about a night market, tastes in food, local flavours, light food and greasy food.
\end{chapteropening}

\section[yèshì]{yèshì — night market}

\label{lesson:ZH-C449-yeshi}



\begin{quote}
Before the new one: say the Mandarin for a market held at dawn, then the Mandarin for to set up a street stall.
\end{quote}

\subsection*{You'll want to know: yèshì}
\emph{yèshì} — \textquotedblleft{}night market\textquotedblright{}.

\begin{quote}
\emph{wǎnshang wǒmen zuò gōngjiāo qù yèshì} — In the evening we take the bus to the night market.
\end{quote}

Its characters come later; for now it is a word for the ear.

\begin{cousinweb}
\emph{yèshì} is built from \emph{yè} 'night' and \zh{市} 'a market', both inherited from Middle Chinese, \emph{yaeH} and \emph{dzyiX} (Guangyun, in Baxter's transcription).
\end{cousinweb}

\begin{grammarlens}[title={What you've built}]
One more word for getting around.
\end{grammarlens}

\subsection*{Your turn}
\begin{itemize}
\raggedright
  \item \emph{Say it:} \emph{yèshì}
  \item \emph{Say it:} \emph{yèshì}, once more
  \item \emph{Say it:} \emph{wǎnshang wǒmen zuò gōngjiāo qù yèshì} — In the evening we take the bus to the night market.
  \item \emph{Recall:} say the Mandarin for a market held at dawn, then the Mandarin for to set up a street stall, then say \emph{yèshì} again
\end{itemize}

\begin{tcolorbox}[breakable,colback=teal!4,colframe=teal!35!black,title={Before you move on}]
What does yèshì mean? (\textquotedblleft{}Night market\textquotedblright{}.) Say it once more.
\end{tcolorbox}
\section[\texorpdfstring{\zh{口味} (kǒuwèi)}{口味 (kǒuwèi)}]{\zh{口味} (kǒuwèi) — one's taste in food}

\label{lesson:ZH-C449-kouwei}



\begin{quote}
Before the new one: say the Mandarin for to set up a street stall, then the Mandarin for a night market.
\end{quote}

\subsection*{You'll want to know: \zh{口味}}
\textbf{\zh{口味}} — \emph{kǒuwèi} — \textquotedblleft{}one's taste in food\textquotedblright{}.

\begin{quote}
\emph{nánfāng rén hé běifāng rén de kǒuwèi bù yíyàng} — People in the south and the north have different tastes in food.
\end{quote}

The new character, in pieces: \zh{味} — eight strokes in two pieces: on the left \zh{口} kǒu, small; on the right a short bar, a longer bar under it, a vertical through both, and a stroke falling to the left and one falling to the right. The next lesson puts it in your hand.

\begin{cousinweb}
\textbf{\zh{口味}} (\emph{kǒuwèi}) is built from \zh{口} 'the mouth' and \zh{味} 'a flavour', both inherited from Middle Chinese, \emph{khuwX} and \emph{mj+jH} (Guangyun, in Baxter's transcription), which glosses the second as the five flavours, sour, salty, sweet, bitter and pungent.
\end{cousinweb}

\begin{grammarlens}[title={What you've built}]
One more word for getting around.
\end{grammarlens}

\subsection*{Your turn}
\begin{itemize}
\raggedright
  \item \emph{Say it:} \emph{kǒuwèi}
  \item \emph{Say it:} \emph{kǒuwèi}, once more
  \item \emph{Say it:} \emph{nánfāng rén hé běifāng rén de kǒuwèi bù yíyàng} — People in the south and the north have different tastes in food.
  \item \emph{Recall:} say the Mandarin for to set up a street stall, then the Mandarin for a night market, then say \emph{kǒuwèi} again
\end{itemize}

\begin{tcolorbox}[breakable,colback=teal!4,colframe=teal!35!black,title={Before you move on}]
What does \zh{口味} mean? (\textquotedblleft{}One's taste in food\textquotedblright{}.) Say it once more.
\end{tcolorbox}
\section[\texorpdfstring{\zh{味} (wèi)}{味 (wèi)}]{\zh{味} — a small \zh{口}, then two bars and a vertical}

\label{lesson:ZH-W449-wei-taste}



\begin{quote}
The word first, then the character inside it.

\begin{itemize}
\raggedright
  \item \emph{Recall:} say \emph{kǒuwèi} — \textbf{R1}, one lesson back, before it decays
\end{itemize}

\begin{quote}
\zh{口} \zh{吃}
\end{quote}

\emph{Read it:} each one aloud, with nothing said first — \emph{kǒu}, \emph{chī}
\end{quote}

\subsection*{Script — a small \zh{口}, then two bars and a vertical}
\hlblockfigure{\detokenize{figures/ZH-W449-wei-taste-filmstrip.pdf}}{How \zh{味} is written, stroke by stroke}

\begin{quote}
\zh{味}
\end{quote}

\textbf{Two pieces, left before right.} Count the pieces, not the strokes.

\textbf{The left piece, \zh{口} kǒu, three strokes}, small and set high.

\textbf{The right piece, five strokes.} A short bar. A longer bar under it. A vertical down through both. Then a stroke falling to the left and a stroke falling to the right, from where the vertical crosses the long bar.

3 + 5 = eight.

\subsection*{Your turn}
\begin{itemize}
\raggedright
  \item \emph{Write it:} \zh{味}, all eight, with the model in front of you
  \item \emph{Cover:} write it again from memory, and count the pieces, not the strokes
\end{itemize}

\begin{tcolorbox}[breakable,colback=teal!4,colframe=teal!35!black,title={Before you move on}]
How many strokes in \zh{味}? (\textbf{Eight}.) Which bar is shorter, the top one or the one under it? (\textbf{The top one}.)

Source: \href{https://raw.githubusercontent.com/chanind/hanzi-writer-data/68d10a4b21150cae5e1ebbd223eed289cf32d90c/data/\%E5\%91\%B3.json}{Hanzi Writer Data \zh{味}, pinned stroke paths and medians}.
\end{tcolorbox}
\section[\texorpdfstring{\zh{风味} (fēngwèi)}{风味 (fēngwèi)}]{\zh{风味} (fēngwèi) — local flavour; distinctive taste}

\label{lesson:ZH-C449-fengwei}



\begin{quote}
Before the new one: say the Mandarin for a night market, then the Mandarin for one's taste in food.
\end{quote}

\subsection*{You'll want to know: \zh{风味}}
\textbf{\zh{风味}} — \emph{fēngwèi} — \textquotedblleft{}local flavour; distinctive taste\textquotedblright{}.

\begin{quote}
\textbf{\zh{我想吃这个岛的风味小吃}} \emph{wǒ xiǎng chī zhège dǎo de fēngwèi xiǎochī} — I'd like to try this island's local snacks.
\end{quote}

\begin{cousinweb}
\textbf{\zh{风味}} (\emph{fēngwèi}) is built from \zh{风} 'a custom' and \zh{味} 'a flavour', both inherited from Middle Chinese, \emph{pjuwng} and \emph{mj+jH} (Guangyun, in Baxter's transcription); the sense 'a local custom' for the first is not given there.
\end{cousinweb}

\begin{grammarlens}[title={What you've built}]
One more word for getting around.
\end{grammarlens}

\subsection*{Your turn}
\begin{itemize}
\raggedright
  \item \emph{Say it:} \emph{fēngwèi}
  \item \emph{Say it:} \emph{fēngwèi}, once more
  \item \emph{Say it:} \emph{wǒ xiǎng chī zhège dǎo de fēngwèi xiǎochī} — I'd like to try this island's local snacks.
  \item \emph{Recall:} say the Mandarin for a night market, then the Mandarin for one's taste in food, then say \emph{fēngwèi} again
\end{itemize}

\begin{tcolorbox}[breakable,colback=teal!4,colframe=teal!35!black,title={Before you move on}]
What does \zh{风味} mean? (\textquotedblleft{}Local flavour; distinctive taste\textquotedblright{}.) Say it once more.
\end{tcolorbox}
\section[qīngdàn]{qīngdàn — light and mild (of food)}

\label{lesson:ZH-C449-qingdan}



\begin{quote}
Before the new one: say the Mandarin for one's taste in food, then the Mandarin for a distinctive local flavour.
\end{quote}

\subsection*{You'll want to know: qīngdàn}
\emph{qīngdàn} — \textquotedblleft{}light and mild (of food)\textquotedblright{}.

\begin{quote}
\emph{zhè jiā fànguǎn de cài hěn qīngdàn} — The food at this restaurant is light.
\end{quote}

Its characters come later; for now it is a word for the ear.

\begin{cousinweb}
\emph{qīngdàn} is built from \emph{qīng} 'clear' and \emph{dàn} 'thin', both inherited from Middle Chinese, \emph{tshjeng} and \emph{damX} (Guangyun, in Baxter's transcription), which glosses the second as a thin taste.
\end{cousinweb}

\begin{grammarlens}[title={What you've built}]
One more word for getting around.
\end{grammarlens}

\subsection*{Your turn}
\begin{itemize}
\raggedright
  \item \emph{Say it:} \emph{qīngdàn}
  \item \emph{Say it:} \emph{qīngdàn}, once more
  \item \emph{Say it:} \emph{zhè jiā fànguǎn de cài hěn qīngdàn} — The food at this restaurant is light.
  \item \emph{Recall:} say the Mandarin for one's taste in food, then the Mandarin for a distinctive local flavour, then say \emph{qīngdàn} again
\end{itemize}

\begin{tcolorbox}[breakable,colback=teal!4,colframe=teal!35!black,title={Before you move on}]
What does qīngdàn mean? (\textquotedblleft{}Light and mild (of food)\textquotedblright{}.) Say it once more.
\end{tcolorbox}
\section[yóunì]{yóunì — greasy; oily}

\label{lesson:ZH-C449-youni}



\begin{quote}
Before the new one: say the Mandarin for a distinctive local flavour, then the Mandarin for light and mild (of food).
\end{quote}

\subsection*{You'll want to know: yóunì}
\emph{yóunì} — \textquotedblleft{}greasy; oily\textquotedblright{}.

\begin{quote}
\emph{yèshì de dōngxi tài yóunì le} — The food at the night market is too greasy.
\end{quote}

Its characters come later; for now it is a word for the ear.

\begin{cousinweb}
\emph{yóunì} is built from \emph{yóu} 'oil' and \emph{nì} 'fatty', both inherited from Middle Chinese, \emph{yuw} and \emph{nrijH} (Guangyun, in Baxter's transcription), which glosses the first as oil and fat and the second as fat and greasy.
\end{cousinweb}

\begin{grammarlens}[title={What you've built}]
One more word for getting around.
\end{grammarlens}

\subsection*{Your turn}
\begin{itemize}
\raggedright
  \item \emph{Say it:} \emph{yóunì}
  \item \emph{Say it:} \emph{yóunì}, once more
  \item \emph{Say it:} \emph{yèshì de dōngxi tài yóunì le} — The food at the night market is too greasy.
  \item \emph{Recall:} say the Mandarin for a distinctive local flavour, then the Mandarin for light and mild (of food), then say \emph{yóunì} again
\end{itemize}

\begin{tcolorbox}[breakable,colback=teal!4,colframe=teal!35!black,title={Before you move on}]
What does yóunì mean? (\textquotedblleft{}Greasy; oily\textquotedblright{}.) Say it once more.
\end{tcolorbox}
